\documentclass[a4paper,12pt]{article}
\usepackage[utf8]{inputenc}
\usepackage[T1]{fontenc}
\usepackage[ngerman]{babel}
\usepackage{amsmath}
\usepackage{amssymb}
\usepackage{amsthm}
\usepackage{enumitem}
\usepackage{geometry}
\geometry{a4paper, margin=2.5cm}

\title{Einsendeaufgaben -- Lektion 1\\[0.5em]
\large Modul 61111: Mathematische Grundlagen}
\author{}
\date{}

\begin{document}
\maketitle

\section*{Aufgabe 1.6}

Sei $A = \mathrm{diag}(a_1, a_2, a_3)$ eine Diagonalmatrix mit paarweise verschiedenen Diagonaleinträgen $a_1, a_2, a_3$ (d.\,h.\ $a_i \neq a_j$ für $i \neq j$).

\textbf{Behauptung:} Jede Matrix $X \in \mathbb{R}^{3 \times 3}$ mit $AX = XA$ ist eine Diagonalmatrix.

\medskip
\textbf{Beweis:}

Sei $X = (x_{ij})_{1 \le i,j \le 3}$. Wir berechnen $AX$ und $XA$:
\[
  AX =
  \begin{pmatrix} a_1 & 0 & 0 \\ 0 & a_2 & 0 \\ 0 & 0 & a_3 \end{pmatrix}
  \begin{pmatrix} x_{11} & x_{12} & x_{13} \\ x_{21} & x_{22} & x_{23} \\ x_{31} & x_{32} & x_{33} \end{pmatrix}
  =
  \begin{pmatrix} a_1 x_{11} & a_1 x_{12} & a_1 x_{13} \\
                  a_2 x_{21} & a_2 x_{22} & a_2 x_{23} \\
                  a_3 x_{31} & a_3 x_{32} & a_3 x_{33} \end{pmatrix},
\]
\[
  XA =
  \begin{pmatrix} x_{11} & x_{12} & x_{13} \\ x_{21} & x_{22} & x_{23} \\ x_{31} & x_{32} & x_{33} \end{pmatrix}
  \begin{pmatrix} a_1 & 0 & 0 \\ 0 & a_2 & 0 \\ 0 & 0 & a_3 \end{pmatrix}
  =
  \begin{pmatrix} a_1 x_{11} & a_2 x_{12} & a_3 x_{13} \\
                  a_1 x_{21} & a_2 x_{22} & a_3 x_{23} \\
                  a_1 x_{31} & a_2 x_{32} & a_3 x_{33} \end{pmatrix}.
\]

Aus $AX = XA$ folgt durch Vergleich der Einträge für $i \neq j$:
\[
  a_i \, x_{ij} = a_j \, x_{ij},
  \quad\text{also}\quad
  (a_i - a_j)\, x_{ij} = 0.
\]
Da $a_i \neq a_j$ für $i \neq j$ vorausgesetzt ist, folgt $x_{ij} = 0$ für alle $i \neq j$.

Damit ist $X$ eine Diagonalmatrix:
\[
  X = \begin{pmatrix} x_{11} & 0 & 0 \\ 0 & x_{22} & 0 \\ 0 & 0 & x_{33} \end{pmatrix}. \qed
\]

\end{document}
